\section{ANTECEDENTES}

\subsection{Recolección de Precios y Monitoreo de Mercado en Bolivia}
Históricamente, el seguimiento de la dinámica de precios al consumidor en el territorio boliviano ha dependido fundamentalmente de relevamientos presenciales coordinados por entidades gubernamentales. El Instituto Nacional de Estadística (INE) fundamenta la estimación del Índice de Precios al Consumidor (IPC) en operativos de campo periódicos en mercados zonales y ferias francas. Por su parte, el Ministerio de Desarrollo Rural y Tierras administra el Observatorio Agroambiental y Productivo (OAP), cuyo objetivo es documentar el comportamiento de precios en origen y en mercados mayoristas para salvaguardar la soberanía y seguridad alimentaria.

El portal oficial del OAP provee un sistema de consulta de precios orientado a productos agrícolas e insumos de la canasta básica. No obstante, su estructura está concebida primordialmente para el consumo visual humano mediante formularios web tradicionales, careciendo de mecanismos de interoperabilidad, servicios web REST o endpoints JSON abiertos que posibiliten el análisis automatizado y continuo de series de tiempo.

\subsection{El Portal de Contrataciones Estatales (SICOES)}
En el ámbito de las compras del Estado boliviano, el Sistema de Contrataciones Estatales (SICOES) es el órgano rector digital establecido bajo el marco del Decreto Supremo N.º 0181 (Normas Básicas del Sistema de Administración de Bienes y Servicios - NB-SABS). Dicho portal concentra los procesos licitatorios de entidades centralizadas, descentralizadas, gobernaciones departamentales y gobiernos autónomos municipales.

A pesar de su trascendencia cívica y económica, el motor de búsqueda del SICOES presenta severas limitaciones de recuperación de información (\textit{Information Retrieval}). El sistema opera bajo concordancias estrictas de subcadenas alfanuméricas sobre el Código Único de Contratación Estatal (CUCE) o el objeto textual declarado. Esto impide la recuperación semántica de convocatorias cuando los términos empleados en los pliegos difieren sintácticamente de la consulta ciudadana (por ejemplo, buscar ``suministro de reactivos de laboratorio'' no retorna licitaciones redactadas como ``adquisición de insumos de diagnóstico clínico''). Asimismo, la plataforma no provee un modelado de redes que vincule las entidades licitantes con sus proveedores recurrentes o consorcios empresariales.

\subsection{Evolución Tecnológica de las Plataformas Web y el Auge de los WAF}
En la última década, las principales corporaciones de comercio retail y farmacéutico de Bolivia modernizaron sus plataformas digitales, transitando desde sitios basados en HTML estático hacia aplicaciones de una sola página (\textit{Single Page Applications} - SPAs) y arquitecturas híbridas renderizadas en servidor (\textit{Server-Side Rendering} - SSR), construidas sobre frameworks como React, Next.js y ecosistemas de e-commerce como Shopify.

Simultáneamente, con el propósito de mitigar ataques distribuidos de denegación de servicio (DDoS) y prevenir la indexación masiva por terceros, las empresas incorporaron cortafuegos de aplicaciones web de última generación (\textit{Next-Generation WAF}), entre los que destaca \textbf{Radware Bot Manager} \citep{radware2024}. Estos sistemas implementan análisis biométrico pasivo, verificación de huella digital de navegador (\textit{browser fingerprinting}) y retos dinámicos de captcha (\textit{hCAPTCHA}), tornando obsoletos a los scrapers tradicionales basados en peticiones HTTP primitivas (como \texttt{curl}, \texttt{requests} o \texttt{cheerio}). Ante este escenario, la investigación en Minería Web demanda la adopción de navegadores automatizados controlados mediante protocolos CDP (\textit{Chrome DevTools Protocol}) y estrategias de persistencia de perfiles de usuario.
